\section{Data}\label{sec:data}

\paragraph{Sources, period, unit.} We use 21 daily market series, January 2007 -- November 2026
(Appendix~\ref{app:data}): 20 from Bloomberg (diesel, Brent, WTI and natural-gas futures for the first two
maturities; implied volatilities of diesel and Brent; copper, gold, US Treasury futures; VIX and 3-month VIX;
S\&P~500 and the US and European energy-sector indices; US and euro-area inflation expectations) and the
ECB EUR/USD reference rate. The unit of observation is a \textbf{trading day} of the NYMEX diesel contract;
after building features (one year of history) and removing the last month, whose outcome is not yet known,
the modelling sample has \textbf{4\,701 days (January 2008 -- October 2026)}.

\paragraph{Diesel price.} The Bloomberg licence available to us did not include the ICE Gasoil contract,
the European diesel benchmark. We use the NYMEX New York Harbor ULSD futures (ticker HO), which prices the
same product (ultra-low-sulphur diesel), is the most liquid diesel contract in the world and is strongly
correlated with Gasoil. Prices in US cents per gallon are converted to \textbf{euros per litre} with the
same-day ECB rate (fixed at 14:15 CET, before the NYMEX settlement). The price varies between 0.15 and
1.29~EUR/l (mean 0.52); Figure~\ref{fig:price} shows the 2008 peak, the 2014--16 and 2020 collapses
and the 2022 energy crisis.

\paragraph{Targets without look-ahead and without roll bias.} A decision taken at the close of day $t$ is
executed at the settlement of $t+1$; the hedged month is days $t+2,\dots,t+22$. The \emph{market cost}
$M_t$ is the average diesel price over these 21 days. The \emph{lockable price} $P_t$ is a strip of futures
settled on $t+1$: each day of the hedged month is priced with the contract that will be front month on that
day (the generic series roll on the first trading day of each month). Comparing different delivery months
would create a spurious gain from locking in backwardated markets: a first version of our analysis made this
mistake and found that always locking saved 4.6\% over 2022--2026, against 2.2\% (before any hedging margin)
with the correct strip. The premium is $R_t=M_t/P_t-1$ and the
spike label is $S_t=\1\{R_t>5\%\}$. Spikes occur on 19.2\% of days, from 0.4\% (2014) to 51\% (2026).

\begin{figure}[t]
\centering\includegraphics[width=\textwidth]{fig1_price_target.pdf}
\caption{Diesel price and target. Top: wholesale diesel price in EUR/litre. Bottom: premium $R_t$ of next
month's market cost over the price lockable at $t$; red areas are spikes ($R_t>5\%$). The shaded period
is the out-of-sample evaluation (2012--2026).}
\label{fig:price}
\end{figure}

\paragraph{Preparation.} (i) Calendars differ across exchanges, so all series are aligned on diesel
trading days and gaps are filled \emph{forward only}, for at most five days. (ii) One implied volatility of
1.7\% (12 September 2007) is a data error; values below 5\% are treated as missing. (iii) WTI was negative on
20 April 2020; it enters only through ratios scaled by the second-month price. (iv) From the 21 raw series
we build \textbf{57 features} in eight groups: returns over 5, 21 and 63 days for eleven markets (33);
futures-curve slopes (4); the diesel crack spread and its change (2); implied and realised volatility,
the volatility risk premium and the VIX term structure (8); inflation expectations and their changes (4);
the price relative to its one-year average (2); the past realised premium and spike rate, lagged so that
they are known at $t$ (2); and month of the year (2). Every feature uses data up to $t$ only. An automated
test scrambles all prices after a date and checks that no feature before that date changes.
